% !TEX root = ../nc-cubic.tex
\section{中心商上の線対象と次数付き持上げの剛性}
\label{mbcr:sec:central-rigidity}

本節では、中心商上で次数一軌道の対象を保つ自己同型は、
元の第三 Veronese 上でもそれらを保つことを示す。
二ブロックで残っていた中心還元核を、自己同値群の有限性を仮定せずに除く。
証明は中心変数が複数あっても同じなので、全ての満支持 $k\geq2$ を扱う。
必要なのは中心上の平坦性、負次数の $\Ext^1$ と $\Hom$ の消滅、
および正次数の中心イデアルによる商の次数別安定性である。

基礎体を標数零の代数閉体 $K$ とし、
第\ref{mbiq:sec:quotients}節の仮定の下で
\[
 B=\bigotimes_{i=1}^kS_i,\qquad
 A=B/\left(\sum_i a_i\Omega_i\right),\quad a_i\ne0,\qquad
 T_i=S_i/(\Omega_i),\qquad E=\bigotimes_iT_i
\]
と置く。本節の $E$ は次数付き代数であり、楕円曲線 $E_i$ と区別する。
さらに
\begin{equation}\label{mbcr:eq:setup}
 C=A^{(3)},\quad Z_1=Z(C)_1,\quad
 R=\Sym_K Z_1,\quad \mathfrak m=R_{>0},\quad I=\mathfrak m C,\quad
 D=C/I=E^{(3)}
\end{equation}
とする。$R$ の生成元の次数は一、$\dim_KZ_1=k-1$ である。
中心一次空間の同定には命題\ref{mbv:prop:center}の $d=1$ を使う。
各整数 $j$ について
\[
 M_j=(A(j))^{(3)},\qquad \overline M_j=M_j/IM_j=(E(j))^{(3)}
\]
と書く。内部次数の規約は $M(j)_a=M_{a+j}$ である。
以下の $\Ext_{\Gr}$ と $\Hom_{\Gr}$ は次数零の射について取る。
内部次数を全て集める場合には下線を付ける。

\subsection{中心上の自由性と完全交叉の局所コホモロジー}

\begin{lemma}[中心多項式環上の自由性]
\label{mbcr:lem:central-free}
$R\to C$ は単射であり、$C$ と全ての $M_j$ は
次数付き $R$-加群として自由である。
自由基底は無限でもよいが、各内部次数は有限次元で、次数は下に有界である。
\end{lemma}
\begin{proof}
$T_i$ の斉次基底を $S_i$ へ持ち上げ、その張る空間を $H_i$ とする。
$\Omega_i$ の正則性と
$S_i/\Omega_iS_i=T_i$ により、乗法は加群同型
\[
 K[z_i]\otimes_KH_i\xrightarrow{\sim}S_i,\qquad z_i\longmapsto\Omega_i,
 \quad\deg z_i=3
\]
を与える。全射性は次数に関する帰納で従う。
線型関係があれば、最小の $z_i$ 冪を正則性によって消去し、
$z_i$ で割った商を取ることで、その最小係数が零となる。
これを繰り返すと単射性も得られる。

従って $H=\bigotimes_iH_i$ とすると、加群として
\[
 B\simeq K[z_1,\ldots,z_k]\otimes_KH,\qquad
 A\simeq
 \left(K[z_1,\ldots,z_k]/\left(\sum_i a_iz_i\right)\right)\otimes_KH .
\]
中心多項式環の次数を三で割ると、括弧内は $R$ となる。
$H$ の次数 $d$ の基底元は、$d\equiv j\pmod3$ のときに限り
$M_j$ の自由基底を与え、その新しい次数は $(d-j)/3$ である。
特に $1\in H_0$ の成分が $R\to C$ の単射性を与える。
剰余を取れば、表示した $\overline M_j$ の同定も従う。
\end{proof}

\begin{lemma}[中心完全交叉の双対と低次切断]
\label{mbcr:lem:ci-cohomology}
有限次元の次数付き右 $E$-加群 $F$ に対し、
係数の一次元空間を自明化して
\begin{equation}\label{mbcr:eq:finite-duality}
 \mathbf R\underline{\Hom}_E(F,E)\simeq F^\vee[-2k]
\end{equation}
を得る。従って、irrelevant torsion の導来関手を
$\mathbf R\Gamma_E$ と書けば
\begin{equation}\label{mbcr:eq:ci-local-cohomology}
 \mathbf R\Gamma_E(E)\simeq E^\vee[-2k].
\end{equation}
特に $E$ は局所コホモロジーが次元 $2k$ に集中する
AS--Cohen--Macaulay 代数である。
$\mathcal O_E=\pi_EE$ とすると、全整数 $b$ について
\begin{equation}\label{mbcr:eq:ci-low-sections}
 \Hom_{\QGr E}(\mathcal O_E,\mathcal O_E(b))=E_b,\qquad
 \Ext^q_{\QGr E}(\mathcal O_E,\mathcal O_E(b))=0
 \quad(1\leq q\leq2k-2)
\end{equation}
となる。
\end{lemma}
\begin{proof}
前補題の tensor 積表示により、$\Omega_1,\ldots,\Omega_k$ は
$B$ 上の中心正則列である。
定理\ref{mbsd:thm:ambient-duality}の有限 Koszul 双対から
\[
 \mathbf R\underline{\Hom}_B(F,B)\simeq F^\vee(3k)[-3k].
\]
一方、この中心正則列の Koszul 解消を双対化すると
\[
 \mathbf R\underline{\Hom}_B(E,B)\simeq E(3k)[-k].
\]
change of rings の導来随伴を適用して二式を比較すると
\eqref{mbcr:eq:finite-duality}を得る。
これは定理\ref{mbsd:thm:gorenstein}の二項解消による計算を
$k$ 回行ったものである。中心性により normalizing twist は生じない。
同じ議論は左右を交換しても成立する。
また任意の $E$-加群に対する change of rings と
$B$ の大域次元 $3k$ から、$E$ の左右単射次元は高々 $2k$ である。
$F=K$ の非零な第 $2k$ Ext により等号となる。

$F_l=E/E_{\geq l}$ は有限次元である。
torsion は
$\Gamma_E(-)=\varinjlim_l\underline{\Hom}_E(F_l,-)$ で計算され、
filtered colimit は完全なので
\[
 R^q\Gamma_E(E)
   =\varinjlim_l\underline{\Ext}^q_E(F_l,E).
\]
\eqref{mbcr:eq:finite-duality}を代入して
\eqref{mbcr:eq:ci-local-cohomology}を得る。
Hilbert 級数は
$(1+t+t^2)^k/(1-t)^{2k}$ なので、ここでの次元は $2k$ である。

切断に移るための局所化三角の入力も確認する。
定理\ref{thm:nfm-relative-line-cohomology}の証明は、
有限 Koszul 解消を持つ $B$ 上で有界幅の derived torsion 模型を構成している。
その中心商への比較は、今回の正則列の各段階で反復できる。
実際、$Q'=Q/(\Omega)$、$\deg\Omega=3$ とし、
$J$ を injective な次数付き右 $Q'$-加群とすると
\[
 \underline{\Ext}^{\,q}_Q(Q/Q_{\geq l},J)
 \simeq
 \underline{\Hom}_{Q'}\!
 \left(\operatorname{Tor}^{Q}_q(Q/Q_{\geq l},Q'),J\right).
\]
正則元の二項解消により $q>1$ の Tor は零である。
第 $1$ Tor は切断の上端にある幅三の次数に限られ、
$l'\geq l+3$ から $l$ への遷移は零となる。
従って忘却した $J$ は $\Gamma_Q$-acyclic であり、
restriction は derived torsion の比較を与える。
これを全段階で反復すれば、$E$ 上でも torsion 複体上で恒等となる
有界幅の derived torsion と Gabriel 局所化三角
\begin{equation}\label{mbcr:eq:localization}
 \mathbf R\Gamma_E(N)\longrightarrow N
 \longrightarrow\mathbf R\omega_E\pi_EN\longrightarrow
\end{equation}
を得る。標準的な局所化三角と切断右随伴の記述については
\cite[\S4, Proposition~4.5 and equation~(4.15)]{NFMBurkeStevenson}も参照されたい。
ここでは上記の有限解消と正則中心列から、その必要な入力を直接確認した。
$N=E$ の内部次数 $b$ に
\eqref{mbcr:eq:ci-local-cohomology}を代入すれば
\eqref{mbcr:eq:ci-low-sections}が従う。
この計算は $\QGr E$ の有限大域次元や smoothness を用いない。
\end{proof}

\subsection{Veronese の根加群に対する負次数 Ext}

\begin{proposition}[根加群の負次数消滅]
\label{mbcr:prop:negative-ext}
任意の整数 $j$ と $n>0$ について
\begin{equation}\label{mbcr:eq:negative-ext}
 \Hom_{\Gr D}(\overline M_j,\overline M_j(-n))=0,\qquad
 \Ext^1_{\Gr D}(\overline M_j,\overline M_j(-n))=0 .
\end{equation}
\end{proposition}
\begin{proof}
補題\ref{lem:nfm-gabriel}の完全同値を
$\mathsf V:\QGr E\xrightarrow{\sim}\QGr D$ と書く。
このとき
\[
 \mathsf V(\mathcal O_E(j))=\pi_D\overline M_j,\qquad
 \mathsf V(\mathcal O_E(-3a))=\mathcal O_D(-a).
\]
切断右随伴を次数 $a$ で評価することは
$\Hom(\mathcal O_D(-a),-)$ に等しい。
従って、右随伴の各次数とその積作用を比較すると、
加群上の Veronese を $\mathsf V_{\Gr}$ として
$\mathsf V_{\Gr}\omega_E\simeq\omega_D\mathsf V$ を得る。
$\mathsf V_{\Gr}$ は完全であり、$\mathsf V$ は完全同値なので injective を保つ。
この関手の同型を導来化すれば、単位射に両立する同型
\begin{equation}\label{mbcr:eq:veronese-sections}
 \mathbf R\omega_D\pi_D\overline M_j
 \simeq
 \left(\mathbf R\omega_E\pi_E E(j)\right)^{(3)}
\end{equation}
を得る。これは切断右随伴の同定であり、
通常の加群 $\Ext$ が Veronese で自動的に保存されるという主張ではない。

式\eqref{mbcr:eq:localization}と
\eqref{mbcr:eq:ci-local-cohomology}により、
$E(j)\to\mathbf R\omega_E\pi_EE(j)$ の cone は
コホモロジー次数 $2k-1$ にのみ成分を持つ。
Veronese は次数付き加群上で完全なので、
\[
 \operatorname{Cone}\!
 \left(\overline M_j\longrightarrow\mathbf R\omega_D\pi_D\overline M_j\right)
 \in D^{\geq2k-1}(\Gr D).
\]
同じことは $\overline M_j(-n)$ にも成立する。
次数零の複体 $\overline M_j$ からの導来 Hom は、この下界を下げない。
従って $q=0,1$ について、導来随伴と上の cone は同型
\begin{align*}
 \Ext^q_{\Gr D}(\overline M_j,\overline M_j(-n))
 &\simeq
 \Ext^q_{\QGr D}(\pi_D\overline M_j,\pi_D\overline M_j(-n))\\
 &\simeq
 \Ext^q_{\QGr E}(\mathcal O_E(j),\mathcal O_E(j-3n))\\
 &\simeq
 \Ext^q_{\QGr E}(\mathcal O_E,\mathcal O_E(-3n))
\end{align*}
を与える。最後の同型は、次数 shift が $\QGr E$ の完全自己同値であることだけを使う。
$\mathcal O_E(j)$ を可換多様体上の線束と仮定していない。
式\eqref{mbcr:eq:ci-low-sections}と $E_{-3n}=0$ が
\eqref{mbcr:eq:negative-ext}を与える。
\end{proof}

\subsection{中心 thickening 上の同型の一意持上げ}

\begin{lemma}[正次数中心に沿った次数付き持上げ]
\label{mbcr:lem:graded-lifting}
$V$ を有限次元 $K$-空間、$R=\Sym_K V$ を $V$ を次数一に置く多項式環とし、
$C$ は中心写像 $R\to C$ を持つ非負次数付き代数で、
$R$ 上平坦であるとする。
$I=R_{>0}C$、$D=C/I$ と置く。
$P,Q$ は下に有界な次数付き右 $C$-加群で、
$R$ 上平坦とし、$\overline P=P/IP$, $\overline Q=Q/IQ$ と書く。
全ての $n>0$ に対して
\[
 \Hom_{\Gr D}(\overline P,\overline Q(-n))
 =\Ext^1_{\Gr D}(\overline P,\overline Q(-n))=0
\]
なら、指定された次数付き同型
$\overline P\xrightarrow{\sim}\overline Q$ は、
ただ一つの次数付き同型 $P\xrightarrow{\sim}Q$ に持ち上がる。
\end{lemma}
\begin{proof}
$\mathfrak m=R_{>0}$、$R_n=R/\mathfrak m^n$、
$C_n=C/I^n$、$P_n=P/I^nP$、$Q_n=Q/I^nQ$ とする。
平坦性により、$C_n$ と $P_n,Q_n$ はそれぞれ $R_n$ 上平坦である。
従って
\begin{equation}\label{mbcr:eq:derived-base-change}
 P_{n+1}\otimes^{\mathbf L}_{C_{n+1}}D
 \simeq
 P_{n+1}\otimes^{\mathbf L}_{C_{n+1}}
       \left(C_{n+1}\otimes^{\mathbf L}_{R_{n+1}}K\right)
 \simeq P_{n+1}\otimes^{\mathbf L}_{R_{n+1}}K
 \simeq\overline P .
\end{equation}
最初の同定には $C_{n+1}$ の平坦性、最後の同定には
$P_{n+1}$ の平坦性を使う。
中心の正則性だけから導来 tensor を通常の tensor に置き換えていない。
導来随伴により、$D$-加群 $L$ について
\[
 \mathbf R\Hom_{\Gr C_{n+1}}(P_{n+1},L)
 \simeq\mathbf R\Hom_{\Gr D}(\overline P,L)
\]
となる。

$Q$ の平坦性は短完全列
\begin{equation}\label{mbcr:eq:thickening-extension}
 0\longrightarrow
 \overline Q(-n)\otimes_K\Sym^nV
 \longrightarrow Q_{n+1}\longrightarrow Q_n\longrightarrow0
\end{equation}
を与える。ここで $\Sym^nV$ は次数零に置いた係数空間であり、
中心単項式の次数 $n$ は $(-n)$ に含めた。
$P_n\simeq Q_n$ が既に与えられたとする。
$P_{n+1}\to P_n\to Q_n$ を持ち上げる障害と持上げの差は、
\eqref{mbcr:eq:thickening-extension}への Hom 完全列により、それぞれ
\[
 \Ext^1_{\Gr C_{n+1}}
  (P_{n+1},\overline Q(-n)\otimes_K\Sym^nV),\qquad
 \Hom_{\Gr C_{n+1}}
  (P_{n+1},\overline Q(-n)\otimes_K\Sym^nV)
\]
に属する。\eqref{mbcr:eq:derived-base-change}と仮定により両者は零である。
従って同型の候補は一意に持ち上がる。

この持上げは同型である。
余核は $\mathfrak m$ で割ると零になり、
$\mathfrak m^{n+1}=0$ により余核そのものが零となる。
全射の核については、$Q_{n+1}$ の $R_{n+1}$-平坦性から
剰余後も左端の完全性が保たれる。
剰余写像が同型なので核も $\mathfrak m$ で割ると零であり、
同じ冪零性により零となる。
これで互いに両立する同型 $P_n\simeq Q_n$ を全ての $n$ で得る。

最後に、形式的同型を実際の次数付き同型へ戻す。
$P,Q$ の次数の下界を固定すると、各内部次数 $a$ について
$n\gg0$ で $(I^nP)_a=(I^nQ)_a=0$ となる。
従ってその次数の同型は実際の $P_a\simeq Q_a$ に安定する。
任意の斉次元による作用の可換性は、始域と終域の両次数で十分大きな
$n$ を取れば有限商上で確認できる。
よって安定した写像は $C$-線型であり、逆も同じ方法で構成される。
一意性も各有限商での一意性から従う。
完備化した加群の存在だけで結論を止めてはいない。
\end{proof}

\subsection{自己同型への適用と中心還元の忠実性}

自己同型の群については、さらに $j_i$ が相異なり $0,1728$ を避けるとする。
位数の仮定は本節冒頭の $p_i,2p_i,3p_i\ne0$ のままでよい。
全 $\mathcal O$-同値群と同定する場合には
定理\ref{mbv:thm:all-o-third}の有限個の捩れ条件も課す。
$\Phi_\alpha$ は右作用を
$m\cdot c=m\alpha^{-1}(c)$ と変更する関手とする。

\begin{theorem}[中心還元からの次数一軌道の認識]
\label{mbcr:thm:reduction-rigidity}
$\alpha\in\Aut_{\mathrm{gr}}(C)$ とし、その中心商上の作用を
$\overline\alpha\in\Aut_{\mathrm{gr}}(D)$ とする。
$j=1,2$ について
\[
 \Phi_{\overline\alpha}(\overline M_j)\simeq\overline M_j
 \quad\text{in }\Gr D
\]
なら
\[
 \Phi_\alpha(M_j)\simeq M_j\quad\text{in }\Gr C .
\]
各剰余同型を指定すると、その持上げは一意である。
従って $\alpha$ の射影的同型類は
$\mathcal K_k=\Aut_{\mathrm{gr}}(A)/K^\times$ に属する。
さらに、中心還元による準同型
\begin{equation}\label{mbcr:eq:faithful-reduction}
 \Aut_{\mathrm{gr}}(C)/K^\times
 \longrightarrow \Aut_{\mathrm{gr}}(D)/K^\times
\end{equation}
は単射である。
これらの結論に $\Aut_{\mathrm{gr}}(C)/K^\times$ の有限性は不要である。
\end{theorem}
\begin{proof}
$\alpha$ は $Z_1$ を保ち、中心多項式環 $R$ の次数付き自己同型を誘導する。
従って補題\ref{mbcr:lem:central-free}により、
$M_j$ と $\Phi_\alpha(M_j)$ はともに $R$ 上平坦で、下に有界である。
後者の剰余は $\Phi_{\overline\alpha}(\overline M_j)$ となる。
指定した剰余同型で二つを同定し、
命題\ref{mbcr:prop:negative-ext}と補題\ref{mbcr:lem:graded-lifting}を適用する。
これで最初の主張を得る。

Veronese 同値で $M_j$ は $\mathcal O_A(j)$ に対応する。
$C$ の次数付き自己同型から来る関手は $(3)$ と可換するので、
$j=1,2$ の対象を保てば全整数の $\mathcal O_A(j)$ を保つ。
定理\ref{mbva:thm:sss}の証明により、その関手は元の次数付き表示の自己同型から来る。
その証明で使うのは点グラフの $p_i,2p_i,3p_i\ne0$ と $j$ の相異性、
および混合二次関係であり、同節の twist 分類用の $6p_i\ne0$ は使わない。
第三 Veronese 上の次数付き自己同型が同じ構造移送を与える核は
grading scalar であるから、$\alpha$ の射影的同型類は $\mathcal K_k$ に属する。

単射性のため、まず $\mathcal K_k$ の $D$ への作用が射影的に忠実であることを確認する。
$h=(h_i)\in\prod_i\mathcal H_i$ の作用が $D_1$ 上のスカラー $\lambda$ だとする。
$i\ne j$ に対し、混合成分
$(T_i)_2\otimes(T_j)_1\subset D_1$ 上では
$(h_i)_2\otimes(h_j)_1=\lambda\operatorname{id}$ である。
二つの可逆写像の tensor 積がスカラーなら各因子もスカラーなので、
全ての $(h_j)_1=\varepsilon_j\operatorname{id}$ となる。
$h_j$ は $\Omega_j$ を固定するため $\varepsilon_j^3=1$ であり、
純成分 $(T_i)_3\subset D_1$ を見ると $\lambda=1$ である。
混合成分の等式は
$\varepsilon_i^2\varepsilon_j=1$、従って $\varepsilon_j=\varepsilon_i$ を与える。
よって $h$ は共通 $\mu_3$ に属し、$\mathcal K_k$ では恒等である。

次に\eqref{mbcr:eq:faithful-reduction}の核の元を取る。
次数スカラーを除くと、$\overline\alpha=\operatorname{id}_D$ となる代表を選べる。
その代表は全ての $\overline M_j$ を保つので、既に示した部分から
$\mathcal K_k$ に属する。上の忠実性によって、この元は恒等である。
\end{proof}

\begin{corollary}[次数一代数への持上げを持つ還元]
\label{mbcr:cor:reduction-lifts}
前定理の $\overline\alpha$ が、$E=\bigotimes_iT_i$ の
次数付き自己同型 $\beta$ の第三 Veronese への制限であるなら、
$\alpha$ の射影的同型類は $\mathcal K_k$ に属する。
\end{corollary}
\begin{proof}
$\Phi_\beta(E(j))\simeq E(j)$ は
$m\mapsto\beta(m)$ による次数付き加群同型である。
第三 Veronese を取ると
$\Phi_{\overline\alpha}(\overline M_j)\simeq\overline M_j$ となるので、
定理\ref{mbcr:thm:reduction-rigidity}を適用する。
\end{proof}

\begin{corollary}[二ブロックの一次中心変形の非消滅]
\label{mbcr:cor:nonzero-class}
$k=2$ とし、$Z(C)_1=Kg$ の生成元を選ぶ。
$C/(g^2)$ が定める次数付き一次変形類は
\[
 \kappa_C\ne0\quad\text{in }HH^2(D,D)_{-1}
\]
である。従って $C$ は中心次数一変数に関して、一次から既に自明でない変形である。
\end{corollary}
\begin{proof}
次の一般的な事実を使う。
標数が二でない体上で、$C$ が連結二次代数、
$g\in Z(C)_1$ が正則、$D=C/(g)$、$D_1\ne0$ とする。
もし $\kappa_C=0$ なら、$D$ 上恒等で
$g\mapsto-g$ となる非自明な射影的自己同型が存在する。

実際、$C/(g^2)$ の積を斉次線型切断で
$\mu_0+g\mu_1$ と書くと、$\mu_1$ は内部次数 $-1$ の
Hochschild cocycle である。
$[\mu_1]=0$ なら、それを与える内部次数 $-1$ の
一 cochain を用いて切断を変更することで、単位と $g$ を保ち、
$D$ 上恒等となる次数付き同型
$D[g]/(g^2)\simeq C/(g^2)$ を得る。
その $D_1$ の像を $C_1$ へ持ち上げ、$C_1=U\oplus Kg$ と書く。
$R_D\subset U\otimes U$ を $D$ の二次関係空間とすると、
各 $r\in R_D$ の $C_2$ における値は $c_rg^2$、$c_r\in K$ となる。
従って $r-c_rg\otimes g$ は $C$ の二次関係である。

これらと $g\otimes u-u\otimes g$ が全二次関係を張る。
実際、関係空間の $U\otimes U$ への射影は $R_D$ へ全射である。
その核に入る関係は
$g\otimes u+v\otimes g+c\,g\otimes g$ と書け、
積が零という条件は $g(u+v+cg)=0$ である。
$g$ の正則性と直和分解から $v=-u$, $c=0$ となる。
よって $U$ を固定し $g$ を $-g$ に送る写像は全二次関係を保ち、
$C$ の次数付き自己同型を与える。
$U\simeq D_1\ne0$ なので、この自己同型は共通 grading scalar ではない。

本節の $C$ は補題\ref{mbv:lem:quadratic}により二次であり、
$g$ の正則性は補題\ref{mbcr:lem:central-free}から従う。
上の自己同型は定理\ref{mbcr:thm:reduction-rigidity}の中心還元の忠実性に反する。
従って $\kappa_C\ne0$ である。
\end{proof}

\begin{remark}[二ブロックと主張の範囲]
$k=2$ では $R=K[g]$、$I=(g)$ であり、
\eqref{mbcr:eq:ci-local-cohomology}の集中次数は四である。
従って中心還元核は有限開に限らず自明である。
一般二次代数に対する高々 $\mu_2$ という有限群の上界や、
一次中心変形類の非消滅をこの結論の仮定として使う必要はない。
後者の非消滅は、逆に系\ref{mbcr:cor:nonzero-class}で結論として得られる。
この系は非消滅を証明するもので、$\kappa_C$ の具体的な Hochschild cocycle や値は計算していない。

ここでの定理は代数閉体上の固定された表示に関する主張である。
$I$-進の非被約商は元の次数付き代数の内部の thickening であり、
任意の非被約パラメータ底上の自己同型層の認識と同一視しない。
また本節だけでは $\Aut_{\mathrm{gr}}(D)$ の全ての元が $E$ へ持ち上がるとは主張しない。
その分類を別途与えれば、系\ref{mbcr:cor:reduction-lifts}が全自己同値群の決定へ接続する。
\end{remark}
